\section{About Color Science}%
\label{sec:about-color-science}

Color science is the scientific domain studying the human perception of color,
its measurement, and characterization. Terminology in a scientific field is
critical to understanding that field and precisely discussing its concepts.
Color is ubiquitous in our lives but typically not well understood.
\ccPar{}
Mark D. Fairchild writes:~\blockcquote[85]{Fairchild2013u}{Almost everyone
knows what color is. After all, they have had firsthand experience of it since
shortly after birth. However, very few can precisely describe their color
experiences or even precisely define color.}
\ccPar{}
The International Commission on Illumination (CIE) defines perceived color as
the~\blockcquote{CIEbk}{characteristic of visual perception that can be
described by attributes of hue\footnote{See sub-section~\ref{subsec:terminology}
for the definitions of hue, brightness, lightness, colorfulness, saturation and
chroma.}, brightness (or lightness) and colourfulness (or saturation or
chroma).}
\ccPar{}
This document is necessarily an overview of color science. The foundational
texts listed below can provide a deeper understanding.
\begin{itemize}
    \item \textit{Digital Color Management: Encoding Solutions, 2nd Edition}
by~\textcite{Madden2007} provides an in-depth analysis of the nature of color
images, digital color encoding, color management systems\footnote{\textit{
Digital Color Management: Encoding Solutions, 2nd Edition} is influential in
modern color management and roots some design decision of ACES.} and digital
color interchange.
    \item \textit{Color Appearance Models}
by~\textcite{Fairchild2013u} is the reference for the fundamental concepts and
phenomena of color appearance.
    \item \textit{The Reproduction of Colour (6th ed.)}
by~\textcite{Hunt2004b} describes the fundamental principles of colorimetry,
color appearance, color reproduction for photography, television, printing and
electronic imaging.
    \item \textit{Digital Video and HD, Second Edition: Algorithms and Interfaces}
by~\textcite{Poynton2012} details the theory and engineering of digital video
systems.
    \item \textit{Color Science: Concepts and Methods, Quantitative Data and Formulae}
by~\textcite{Wyszecki2000bi} is the authoritative reference for colorimetry.
\end{itemize}

The authors of this document are greatly indebted to the writers of these
books.
\ccPar{}
Dr.~Robert W.G. Hunt passed away during the writing of this section. His work
is foundational to the color science community. Dr.~Hunt’s contributions to the
understanding of human color perception and its application to the engineering
of photographic systems laid the foundation for much of the color science that
is used in the motion picture and gaming industries today. Besides being a
brilliant scientist, Dr.~Hunt was a kind man always willing to share his
knowledge, whether formally in one of the many classes he taught, or informally
over coffee or a meal. Dr.~Hunt’s true passion was teaching others. A few of
the authors of this document had the good fortune to have known and learned
from Dr.~Hunt and were saddened to hear of his passing. His contributions to
the field, friendship and mentorship will be missed.

\begin{figure}[H]
    \includegraphics[alt={Hunt - Tribute.},width=\textwidth]{hunt-tribute.jpg}
    \caption{
        A tribute to Dr.~Robert W.G. Hunt: Light dispersion through a prism
        (bottom right) with the rainbow of color components lighting up
        \textit{The Reproduction of Colour (6th ed.)} by~\textcite{Hunt2004b}.
    }%
    \label{fig:hunt-tribute}
\end{figure}

Color management depends fundamentally on colorimetry, the measurement of
color. Without colorimetry, it is not possible to characterize cameras and
displays or understand the imaging principles that permeate the rest of
computer graphics. While it is possible to immediately dive into color
pipelines, having a rudimentary understanding of concepts such as spectral
measurement, tristimulus values, or color appearance provides a richer
understanding of why some approaches to color management are successful, and
some are not.

This section endeavors to stay generic and abstract. It does not focus on
particular workflow details but provides the necessary foundations for
understanding section 3. This section abundantly references the history of
color science because the understanding of the past is essential for the
comprehension of the present: knowing how colorimetry evolved explains the
current practices. The physical nature of light and the quantities used for
its analysis are examined. The anatomy of the human visual system (HVS) is
reviewed with a particular focus on the eye and the primary adaptation
mechanisms. Basic colorimetry is presented, providing the requisites for
converting spectral measurements to RGB values or understanding perceptual
uniformity. Advanced colorimetry, the gateway to color appearance modeling, is
summarily exposed. Then, the representation of color is described, driving the
section toward concrete concepts in direct relation with motion picture color
management. Finally, color imaging systems and the image capture, signal
processing, and image formation stages are explained using digital imaging
devices such as motion pictures cameras or LCD displays.

\begin{ccKeyPoints}
    \begin{itemize}
        \item Color science studies the human perception of color,
its measurement, and characterization.
        \item Proper terminology usage is critical to the understanding of a
scientific field.
        \item Color is the characteristic of a visual perception.
        \item Color management depends on colorimetry, the measurement of color.
    \end{itemize}
\end{ccKeyPoints}
